As shown in the previous section, R-GCNs serve as an effective encoder for relational data. We now combine our encoder model with a scoring function (which we will refer to as a decoder, see Figure \ref{fig:model-c}) to score candidate triples for link prediction in knowledge bases.

\subsubsection{Datasets}
Link prediction algorithms are commonly evaluated on FB15k, a subset of the relational database Freebase, and WN18, a subset of WordNet containing lexical relations between words. In \citet{toutanova2015observed}, a serious flaw was observed in both datasets: The presence of inverse triplet pairs $t = (e_1, r, e_2)$ and $t' = (e_2, r^{-1}, e_1)$ with $t$ in the training set and $t'$ in the test set. This reduces a large part of the prediction task to memorization of affected triplet pairs. A simple baseline LinkFeat employing a linear classifier on top of sparse feature vectors of observed training relations was shown to outperform existing systems by a large margin. To address this issue, Toutanova and Chen proposed a reduced dataset FB15k-237 with all such inverse triplet pairs removed. We therefore choose FB15k-237 as our primary evaluation dataset. Since FB15k and WN18 are still widely used, we also include results on these datasets using the splits introduced by \citet{bordes2013translating}.

\begin{table}[ht]
\centering
\begin{tabular}{lrrr}
\toprule
Dataset & WN18 & FB15K & FB15k-237\\ \midrule
Entities    & 40,943  & 14,951 & 14,541\\
Relations   & 18 & 1,345 & 237\\
Train edges & 141,442 & 483,142 & 272,115\\
Val.~edges & 5,000 & 50,000 & 17,535\\
Test edges & 5,000 & 59,071 & 20,466\\ \bottomrule
\end{tabular}
\caption{Number of entities and relation types along with the number of edges per split for the three datasets.\label{table:datasets}}
\end{table}

\paragraph{Baselines}
A common baseline for both experiments is direct optimization of \textit{DistMult}~\cite{distmult-embedding_entities_and_relations}. This factorization strategy is known to perform well on standard datasets, and furthermore corresponds to a version of our model with fixed entity embeddings in place of the R-GCN encoder as described in Section \ref{section:link_prediction}. As a second baseline, we add the simple neighbor-based LinkFeat algorithm proposed in \citet{toutanova2015observed}.

We further compare to ComplEx \cite{complex-complex_embeddings_for_simple_link_prediction} and HolE \cite{nickel2015holographic}, two state-of-the-art link prediction models for FB15k and WN18. ComplEx facilitates modeling of asymmetric relations by generalizing DistMult to the complex domain, while HolE replaces the vector-matrix product with circular correlation. Finally, we include comparisons with two classic algorithms -- CP~\cite{hitchcock1927expression} and TransE~\cite{bordes2013translating}.

\begin{table*}[t!]
\centering
\begin{tabular}{@{}lllllllllll@{}}
\toprule
      & \multicolumn{5}{c}{FB15k}                             & \multicolumn{5}{c}{WN18}                            \\ \cmidrule(lr){2-6} \cmidrule(l){7-11}
      & \multicolumn{2}{c}{MRR} & \multicolumn{3}{c}{Hits @} & \multicolumn{2}{c}{MRR} & \multicolumn{3}{c}{Hits @} \\ \cmidrule(lr){2-3} \cmidrule(lr){4-6} \cmidrule(lr){7-8} \cmidrule(l){9-11}
Model & Raw      & Filtered     & 1       & 3      & 10      & Raw      & Filtered     & 1       & 3      & 10      \\ \midrule
LinkFeat 		&  & 0.779 & & & 0.804 & & 0.938 & & & 0.939 \\
\midrule
DistMult 		& 0.248 & 0.634 & 0.522 & 0.718 & 0.814 & 0.526 & 0.813 & 0.701 & 0.921 & 0.943 \\
R-GCN  			& 0.251 & 0.651 & 0.541 & 0.736 & 0.825 & 0.553 & 0.814 & 0.686 & 0.928 & 0.955 \\
R-GCN+			& \textbf{0.262} & \textbf{0.696} & \textbf{0.601} & \textbf{0.760} & \textbf{0.842} & 0.561 & 0.819 & 0.697 & 0.929 & \textbf{0.964}  \\ \midrule
CP* & 0.152 & 0.326 & 0.219 & 0.376 & 0.532 & 0.075 & 0.058 & 0.049 & 0.080 & 0.125 \\
TransE* 		& 0.221 & 0.380 & 0.231 & 0.472 & 0.641 & 0.335 & 0.454 & 0.089 & 0.823 & 0.934 \\
HolE** 			& 0.232 & 0.524 & 0.402 & 0.613 & 0.739 & \textbf{0.616} & 0.938 & 0.930 & \textbf{0.945} & 0.949\\
ComplEx* 		& 0.242 & 0.692 & 0.599 & 0.759 & 0.840 & 0.587 & \textbf{0.941} & \textbf{0.936} & \textbf{0.945} & 0.947 \\ \bottomrule
\end{tabular}
\caption{Results on the the Freebase and WordNet datasets. Results marked (*) taken from \citet{complex-complex_embeddings_for_simple_link_prediction}. Results marks (**) taken from \citet{nickel2015holographic}. R-GCN+ denotes an ensemble between R-GCN and DistMult -- see main text for details.}
\label{table:link-results}
\end{table*}

\paragraph{Results}
We provide results using two commonly used evaluation metrics: \textit{mean reciprocal rank} (MRR) and \textit{Hits at $n$} (H@n).
Following \citet{bordes2013translating}, both metrics can be computed in a raw and a filtered setting. We report both filtered and raw MRR (with filtered MRR typically considered more reliable), and filtered Hits at 1, 3, and 10.

We evaluate hyperparameter choices on the respective validation splits. We found a normalization constant defined as $c_{i,r}=c_{i}=\sum_{r} |\mathcal{N}^r_i|$ --- in other words, applied across relation types -- to work best. For FB15k and WN18, we report results using basis decomposition (Eq. \ref{eq:basis}) with two basis functions, and a single encoding layer with $200$-dimensional embeddings. For FB15k-237, we found block decomposition (Eq. \ref{eq:block}) to perform best, using two layers with block dimension $5 \times 5$ and $500$-dimensional embeddings. We regularize the encoder through edge dropout applied before normalization, with dropout rate $0.2$ for self-loops and $0.4$ for other edges. Using edge droupout makes our training objective similar to that of denoising autoencoders~\cite{vincent2008}. We apply $l2$ regularization to the decoder with a penalty of $0.01$.

We use the Adam optimizer \cite{kingma2014adam} with a learning rate of $0.01$. For the baseline and the other factorizations, we found the parameters from \citet{complex-complex_embeddings_for_simple_link_prediction} -- apart from the dimensionality on FB15k-237 -- to work best, though to make the systems comparable we maintain the same number of negative samples (i.e.~$\omega=1$). We use full-batch optimization for both the baselines and our model.

\begin{figure}
\centering
\begin{tikzpicture}
\begin{axis}[xlabel=Average degree,
ylabel=MRR,
legend style={font=\fontsize{9}{11}\selectfont},
width=1.2\linewidth,
height=0.6\linewidth,
scale=.65]
\addplot[color=red,mark=triangle*] table[x=Degree,y=GCN] {data/degree.dat};
\label{p1}
\addplot[color=blue,mark=x] table[x=Degree, y=Distmult] {data/degree.dat};
\label{p2}
\legend{R-GCN, DistMult}
\end{axis}%
\end{tikzpicture}\vspace{-0.5em}
\caption{Mean reciprocal rank (MRR) for R-GCN  and DistMult on the FB15k validation data as a function of the node degree (average of subject and object).\label{figure:degree}}
\end{figure}

On FB15k, local context in the form of inverse relations is expected to dominate the performance of the factorizations, contrasting with the design of the R-GCN model. To better understand the difference, we plot in Figure \ref{figure:degree} the FB15k performance of the best R-GCN model and the baseline (DistMult) as functions of degree of nodes corresponding to entities in the considered triple
(namely, the average of degrees for the subject and object entities). It can be seen that our model performs better for nodes with high degree where contextual information is abundant. The observation that the two models are complementary suggests combining the strengths of both into a single model, which we refer to as R-GCN+. On FB15k and WN18 where local and long-distance information can both provide strong solutions, we expect R-GCN+ to outperform each individual model. On FB15k-237 where local information is less salient, we do not expect the combination model to outperform a pure R-GCN model significantly. To test this, we evaluate an ensemble (R-GCN+) with a trained R-GCN model and a separately trained DistMult factorization model: $f(s,r,t)_{\text{R-GCN+}} = \alpha f(s,r,t)_{\text{R-GCN}} + (1- \alpha) f(s,r,t)_{\text{DistMult}}$, with $\alpha=0.4$ selected on FB15k development data.

In Table~\ref{table:link-results}, we evaluate the R-GCN model and the combination model (R-GCN+) on FB15k and WN18.

On the FB15k and WN18 datasets, R-GCN and R-GCN+ both outperform the DistMult baseline, but like all other systems underperform on these two datasets compared to the LinkFeat algorithm. The strong result from this baseline highlights the contribution of inverse relation pairs to high-performance solutions on these datasets. Interestingly, R-GCN+ yields better performance than ComplEx for FB15k, even though the R-GCN decoder (DistMult) does not explicitly model asymmetry in relations, as opposed to ComplEx.

This suggests that combining the R-GCN encoder with the ComplEx scoring function (decoder) may be a promising direction for future work. The choice of scoring function is orthogonal to the choice of encoder; in principle, any scoring function or factorization model could be incorporated as a decoder in our auto-encoder framework.

\begin{table}[htp!]
\centering
\begin{tabular}{llllll}
\toprule
                                  & \multicolumn{2}{c}{MRR} & \multicolumn{3}{c}{Hits @} \\ \cmidrule{2-3} \cmidrule{4-6}
Model                             & Raw      & Filtered     & 1       & 3       & 10     \\ \midrule
LinkFeat     &  & 0.063         &    & & 0.079  \\ \midrule
DistMult     & 0.100         & 0.191         & 0.106   & 0.207   & 0.376  \\
R-GCN               & \textbf{0.158}       & 0.248 & \textbf{0.153}     & 0.258     & 0.414     \\
R-GCN+                     & 0.156 & \textbf{0.249} & 0.151     & \textbf{0.264}     & \textbf{0.417} \\     \midrule
CP     & 0.080 & 0.182 & 0.101 & 0.197 & 0.357 \\
TransE    	& 0.144	& 0.233 & 0.147 & 0.263 &	0.398  \\
HolE    	& 0.124	& 0.222 & 0.133 & 0.253 &	0.391  \\
ComplEx    & 0.109 & 0.201 & 0.112 & 0.213   & 0.388 \\
\bottomrule
\end{tabular}
\caption{Results on FB15k-237, a reduced version of FB15k with problematic inverse relation pairs removed. CP, TransE, and ComplEx were evaluated using the code published for \citet{complex-complex_embeddings_for_simple_link_prediction}, while HolE was evaluated using the code published for \citet{nickel2015holographic}.}
\label{table:fb15k-237}
\end{table}

In Table \ref{table:fb15k-237}, we show results for FB15k-237 where (as previously discussed) inverse relation pairs have been removed and the LinkFeat baseline fails to generalize\footnote{Our numbers are not directly comparable to those reported in \citet{toutanova2015observed}, as they use pruning both for training and testing (see their sections 3.3.1 and 4.2). Since their pruning schema is not fully specified (e.g., values of the relation-specific parameter $t$ are not given) and the code is not available, it is not possible to replicate their set-up.}. Here, our R-GCN model outperforms the DistMult baseline by a large margin of 29.8\%, highlighting the importance of a separate encoder model. As expected from our earlier analysis, R-GCN and R-GCN+ show similar performance on this dataset. The R-GCN model further compares favorably against other factorization methods, despite relying on a DistMult decoder which shows comparatively weak performance when used without an encoder.
